
\paragraph{\bf{Holographic Near-eye Displays.}} Holographic displays are a promising technology for virtual and augmented reality applications due to their unique capability to display true 3D content and significant progress has been made recently~\cite{Javidi2021-vh, Chang2020-ja, Pi2022-mr}. In particular, the advancement in computer graphics, machine learning, and computing infrastructures have enabled real-time hologram rendering based on neural networks~\cite{shi2021neural, peng2020neural}, significantly improved image quality with end-to-end optimization~\cite{peng2020neural,chakravarthula2020learned,choi2022flex}, higher light efficiency and brightness with simultaneous control of multiple wavelengths and energy-efficiency loss function~\cite{Markley2023-eg, Chao2023-wl, Kavakli2023-nd}, and thin form factors in eyeglasses-like design~\cite{maimone2017holo,jang2024waveguide, Kim2022-cd,gopakumar2024ar}. Despite offering these unique capabilities, current holographic displays fail to provide a comfortable immersive experience as they cannot simultaneously provide a wide field of view (FoV) and a sufficiently large eyebox (i.e., the region in which a user's eye perceives the displayed content). 

In a given display system, the product of the FoV and the eyebox is a constant, referred to as the {\it \'etendue}. For a holographic display, the \'etendue is directly proportional to the number of pixels in the SLM. A 1080p SLM, for instance, can either support a wide field of view (e.g., 80 degrees) with an eyebox smaller than 1~mm or vice versa. %To achieve a comfortable and immersive experience, an SLM with a resolution surpassing 10k is essential. 
However, increasing the SLM resolution to the point where both large field of view and eyebox can be achieved simultaneously faces significant challenges in manufacturing, cost, and addressing speed, accuracy, and bandwidth. Instead, efforts have been made to increase the \'etendue of holographic displays without increasing SLM resolution.
The approaches fall under two categories: (i) increasing \'etendue {\it after the SLM} and (ii) increasing \'etendue {\it before} the SLM.


% 20240214 talk about anat's paper
\paragraph{\bf{Post-SLM \'Etendue Expansion}}
The most representative method in this category is mask-based \'etendue expansion \cite{kuo2020mask, tseng2023neural, park2019photon, yu2017speckle, buckley2006mask, monin2022mask}, where a static mask at a resolution higher than the SLM is placed after the SLM to increase the diffraction angles of the SLM-modulated wavefront, thus increasing the \'etendue. However, such systems suffer from difficulties in alignment and from reduced image quality and low contrast because their effective degrees of freedom~\cite{starikov1982effective} are insufficient to synthesize a high-quality large-\'etendue wavefront. Pupil replication \cite{bigler2018expansion, bigler2019expansion, draper2019expansion, draper2021dof, draper2022expansion, jang2024waveguide, Park2018-ge,kress2020waveguide} is another popular approach. It is implemented either by putting a pupil-replicating waveguide after the SLM to replicate pupil locations at its out-coupler or using a pupil-replicating holographic optical element (HOE) as the eyepiece, effectively expanding the eyebox of the system. However, pupil-replicating displays cannot display 3D content or natural parallax across the expanded eyebox since the content within the eyebox are merely copies of the same wavefront. Higher-diffraction orders combined with pupil optimization can also be leveraged to slightly expand the eyebox in the single-source case \cite{shi2024ergo,schiffers2023stochastic}. Finally, a regular eyepiece can be replaced by a lens array to partition an unexpanded eyebox into an array of smaller chunks that cover an expanded area~\cite{Chae2023-ll, Wang2023-vf, Xia2020-sa}. However, this comes at an explicit cost of image quality and brightness nonuniformity, especially when observed with a small pupil.

\paragraph{\bf{Pre-SLM \'Etendue Expansion}}
Methods in this category modify the laser illumination to expand \'etendue either through a multi-source configuration or beam-steering. \citet{jang2018hoe} used a micro-electromechanical-system (MEMS) mirror to temporally change the laser illumination and steer the resulting pupils over a larger eye box. \citet{lee2020wide} implemented the same principle by arranging individual laser diodes into a 2D array and sequentially turning each one on to create temporal directional illumination. \citet{monin2022tilt}  implemented \textit{per-pixel} beam steering of the phase SLM by using transmissive LCD panels and polarization gratings and demonstrated that the \'etendue expansion amount scales exponentially with the number of LCD layers.
To permanently expand the \'etendue, \citet{jo2022binary} activated all illumination sources simultaneously. They introduced a random mask at the Fourier plane to break the correlation among copies in the spectrum formed by directional illuminations. This effectively eliminates duplicate images within the expanded eyebox. However, they did not demonstrate view-dependent effects across the expanded eyebox and the random mask is not content adaptive, resulting in reduced 3D realism and low image quality. Instead of using multiple laser diodes that are incoherent with each other, \citet{lee2022envelope} implemented a mutually coherent multi-laser source using a lens array. The mutually coherent sources can interfere constructively and destructively with each other, granting the hologram optimization process more degrees of freedom. However, their system requires eye tracking and a new hologram needs to be optimized for each dynamic pupil location, making the system challenging for real-time applications. 

\paragraph{\bf{Fourier Modulation.}} 
Holographic display systems often require Fourier plane filtering to remove HDOs created by the pixelated structure of the phase SLM \cite{maimone2017holo, peng2020neural, shi2021neural, Shi2022-ih}. However, it is not straightforward to apply this to a multi-source or beam steering setting since the directional illuminations create shifted copies of the wavefront from normally incident illumination and the associated HDOs in the Fourier domain. When using beam-steering, the filter position needs to be dynamically adjusted to block the HDOs of the shifted wavefront. To achieve this, \citet{lee2020wide} placed a programmable polarization shutter at the Fourier plane and synchronized the laser sources to filter out the HDOs. However, when using multi-source illumination for \'etendue expansion, HDOs associated with one illumination intermingle with the shifted wavefront of another illumination at the Fourier domain, making it impossible to separate. It is therefore crucial to model HDOs to precisely characterize how they contribute to different angular views. Explicit modeling of HDOs has been demonstrated for 2D images~\cite{gopakumar2021unfiltered} and 3D focal stacks~\cite{Kim2022-cd, shi2024ergo}, but not for 4D light fields under multi-source illumination.

\bigskip

Inspired by \citet{jo2022binary}, we employ a multisource laser illumination and additionally place a dynamic amplitude SLM at the Fourier plane to enable content adaptive modulation and time multiplexing. We jointly optimize the patterns for the phase and amplitude SLMs to reproduce a 4D light field rendered over an expanded eyebox. We also explicitly model the HDOs and demonstrate notable improvement in image quality and contrast. Collectively, this new hardware and software co-design enables a dynamic view-dependent holographic display with a large eyebox.



